\section{综合结论与改进方向}\label{sec:conclusions}
\subsection{研究闭环的价值体现在状态连续性}
本轮实验展示了从初始程序到独立验收产物的完整过程。搬箱案例中的升级、算法改写、退步和恢复都发生在同一个任务状态中；c5 的得分退步通过外部验证被识别，c6 继续使用 c4 的实现作为起点。最佳候选保留机制将一次有价值的改进保存下来，为后续实验提供稳定的起点。

固定模型实验的 15 个交付全部通过最终测试，其中 8 个保留初始程序。候选通过率为 62.5\%，反映搜索过程仍包含较多无效尝试。两组数字共同描述了系统行为：研究可以经历失败，最终交付仍由已经通过验证的产物承担。

\subsection{反馈的有效性取决于它怎样改变下一次实验}
ahc024 的失败集中在约束满足；反馈组的一半请求用于分析，却没有得到高于初始程序的合法结果。ahc026 反馈组的 20 个候选只有 9 种精确代码，14 个正确候选中仅 2 个刷新当时最佳。可用反馈已经进入上下文，但尚未持续转化为新的有效操作。

这使下一步优化具有明确的落点：将裁判错误归纳为可执行检查；在提出候选前识别相同代码；为每次提案记录它要验证的具体变化。对于地图问题，可围绕保持连通性与邻接的局部操作组织搜索；对于搬箱问题，可延续已有有效策略，比较目标堆选择和搜索密度等具体变量。

\subsection{动作、模型和停止需要共同分配预算}
路由实验中，动作 Decisions 组生成 7 个正确候选，并在规则升级后获得最高得分。双 Decisions 全程使用 Qwen，在较早阶段结束，以更少时间获得较小幅度的改进。动作决策控制实验类型，模型决策控制生成能力与延迟，停止决策决定是否继续支付下一轮成本；三者共同影响最终结果。

这组结果为下一轮实验提供了可检验的问题：在相同质量门槛下，哪些任务值得提前升级；一次没有刷新最佳之后，应继续改写、改用模型还是停止；路由调用的频率如何随剩余请求下降。扩展到多个任务并重复四种组合，可以进一步估计这些决策的稳定收益。

\subsection{算法层面的解释让结果更有用途}
分数和倍率可以进一步映射到实际工作。ahc026 的最佳方案把 15 例搬运能耗减少了 44.707\%；卷积通过变换路径与受控精度选择减少计算和内存开销；PageRank 通过直接使用稀疏结构缩短执行路径。这些解释把“模型生成了更快或更高分的代码”展开为具体可研究的算法变化。

本轮成果由此形成三层内容：可执行的研究与调度系统、带有独立评测和资源账本的实验记录、能够追踪到候选变化的分析结论。后续迭代可以围绕这些已观察到的现象建立假设，保持任务、预算和评测接口一致，再检验实现改动的效果。
